\note{sa page 3}

\struck{\ill{}} et \struck{que} l'on suppose que

\[ (3.0) \qquad \pi_{i} = 0 \ \text{ si } i \notin [0, 2n], \qquad n \ \text{entier \add{fixé}} . \]

\struck{Supp} On s'occupe ici de l'unicité. \struck{Supposons établi
l'unicité} Notons d'abord que
\[ u \in \mathcal{E}^{0} \iff \bigl( u \pi_{i} = \pi_{i} u \ \ \forall i \bigr) , \]
donc $\mathcal{E}^{0}$ est le commutant de la $k$-algèbre engendrée par les
$\pi_{i}$. Donc les $\pi_{i}$ sont dans le centre
$\mathfrak{z}\mathcal{E}^{0}$ de $\mathcal{E}^{0}$, et ce sont les projecteurs
\add{mutuellement} orthogonaux de somme $1$ de cet anneau commutatif.

Supposons \struck{déterminés} déjà les $\pi_{\alpha}$ pour
$\alpha \notin [i, 2n-i]$ \add{$i < n$}, et déterminons $\pi_{i}$ et
$\pi_{2n-i}$. Posons
\[ \varphi_{i} = 1 - \sum_{\alpha < i} \pi_{\alpha} = \sum_{\alpha \geqslant i} \pi_{\alpha} , \qquad \psi_{2n-i} = 1 - \sum_{\alpha > 2n-i} \pi_{\alpha} = \sum_{\alpha \leqslant 2n-i} \pi_{\alpha} , \]
donc aussi les composés
\[ \psi_{2n-i}\, \mathcal{E}^{2n-2i}\, \varphi_{i} = \sum_{\alpha} \psi_{2n-i}\, \pi_{\alpha + (2n-2i)}\, \mathcal{E}\, \pi_{\alpha}\, \varphi_{i} . \]

\marginal{C'est clair si $i = n-1$, car on connaît déjà les $\pi_{\alpha}$ pour
$\alpha \neq n$, d'où aussi $\pi_{n} = 1 - \sum_{\alpha \neq n} \pi_{\alpha}$.
Mais pour $i \leqslant n-1$ \ill{}}

Or
\[ \pi_{\alpha}\, \varphi_{i} = \begin{cases} 0 & \text{si } \alpha < i \\ \pi_{\alpha} & \text{si } \alpha \geqslant i \end{cases} \qquad \text{et} \qquad \psi_{2n-i}\, \pi_{\alpha+(2n-2i)} = \begin{cases} 0 & \text{si } \alpha > i \\ \pi_{\alpha+(2n-2i)} & \text{si } \alpha \leqslant i \end{cases} \]
donc le terme d'indice $\alpha$ dans la somme est nul sauf si $\alpha = i$,
auquel cas il est égal à
\[ \pi_{i+(2n-2i)}\, \mathcal{E}\, \pi_{i} = \pi_{2n-i}\, \mathcal{E}\, \pi_{i} = \mathcal{E}_{2n-i,\,i} . \]

Donc la connaissance des $\pi_{\alpha}$ pour $\alpha \notin [i, 2n-i]$ implique
la connaissance de $\mathcal{E}_{2n-i,\,i}$ \struck{considéré}. La condition qui
va impliquer l'unicité de $u_{i}$,

